\section{Open contracts and proof-graph compilation}
\label{sec:contracts}

The purpose of a contract is to make a local theorem independent of one
preassembled ambient context.  DARC uses a definition that is exactly
equivalent to an ordinary sequent, so contracts are an organizational layer,
not a new logic.

\subsection{Open contracts}

For a list \(A=[A_1,\ldots,A_k]\) and formula \(\phi\) in the same finite
scope, define in the metatheory
\[
  \Contract(A,\phi)
  \;:\!\!\Longleftrightarrow\;
  \forall\Delta\,
  \bigl((\forall i\le k\;\Delta\vdash A_i)
  \to \Delta\vdash\phi\bigr).
\]

\begin{theorem}[Contract/sequent equivalence]
\label{thm:contract-sequent}
\[
  \Contract(A,\phi)
  \quad\Longleftrightarrow\quad
  A\vdash_{\PRAAx}\phi.
\]
\end{theorem}

\begin{proof}
If the contract holds, choose \(\Delta=A\).  Each \(A_i\) is derivable by the
hypothesis rule, so the contract yields \(A\vdash\phi\).  Conversely, given
\(A\vdash\phi\) and derivations \(\Delta\vdash A_i\), apply simultaneous cut
\cref{lem:simcut}.  This is the proof implemented in
\pathcode{lean/DARC/StructuralKernel.lean}.
\end{proof}

The contract form is convenient for a compiler because it describes a
function from supplied dependencies to a conclusion.  The sequent form is the
canonical proof object checked by Lean.

\subsection{Substitution and context independence}

A local theorem may be proved once with formal parameters
\(x_1,\ldots,x_n\).  At a call site, the compiler substitutes arbitrary terms
\(t_1,\ldots,t_n\), including compound terms produced by other nodes.  By the
checked substitution theorem,
\[
  A[\vec t/\vec x]\vdash\phi[\vec t/\vec x].
\]
It may then weaken the context or use the contract interpretation to insert
the theorem under unrelated local hypotheses.  Thus a node interface should
avoid closing parameters unless closure is part of the mathematical
statement.

\subsection{Acyclic schedule compilation}

Assume the retained proof graph is acyclic and let
\(v_1,\ldots,v_N\) be a topological order.  Each node has a local theorem
\[
  \bigl(\phi_u:u\in\operatorname{dep}(v_i)\bigr),\;B_i
  \vdash \phi_{v_i},
  \tag{7.1}\label{eq:local-contract}
\]
where every formula in \(B_i\) occurs in the global residual context \(B\).

\begin{theorem}[Acyclic graph compilation]
\label{thm:acyclic-compile}
For every retained node \(v_i\), \(B\vdash\phi_{v_i}\).  In particular, if
\(q\) is the output node, then \(B\vdash\phi_q\).
\end{theorem}

\begin{proof}
Induct on the topological order.  Every dependency of \(v_i\) occurs earlier,
so the induction hypothesis gives \(B\vdash\phi_u\) for each dependency.
Every formula in \(B_i\) is a member of \(B\) and is therefore available by
the hypothesis rule.  Apply simultaneous cut to the local theorem
\eqref{eq:local-contract}.  This yields \(B\vdash\phi_{v_i}\).
\end{proof}

The proof is constructive and directly implementable as recursion over the
finite schedule.  The compiler need not add a primitive cut constructor.

\subsection{Multiple outputs and shared subproofs}

If a query asks for several formulas \(\phi_{q_1},\ldots,\phi_{q_s}\), seed all
output coordinates at once.  Demand closure computes the union of their least
backward slices, including interactions through precision and guards.  Compile
each retained node once, then combine the outputs by conjunction introduction
or leave them as a family of sequents.  This avoids duplicating expensive local
proofs.

A theorem library can expose a node with several observations but separate
activation coordinates.  For example, a Euclidean-algorithm node may provide
only a gcd identity for one query, and additionally expose Bezout coefficients
for another.  The latter query activates certificate fields and arithmetic
lemmas not needed by the former.

\subsection{Residual boundary discharge}

Let the compiled endpoint be
\[
  B_1,\ldots,B_r,\Gamma\vdash\Phi.
\]
Suppose later developments prove a subset \(B_i\) under \(\Gamma\).  Repeated
or simultaneous cut removes exactly those boundaries.

\begin{corollary}[Incremental producer discharge]
\label{cor:incremental-discharge}
For \(S\subseteq\{1,\ldots,r\}\), if
\(\Gamma\vdash B_i\) for every \(i\in S\), then
\[
  (B_j:j\notin S),\Gamma\vdash\Phi.
\]
\end{corollary}

This permits useful partial formalization.  A difficult analytic theorem can
be compiled down to a small list of quantitative producer lemmas.  Each
producer can be proved independently without changing the downstream proof.

\subsection{Local theorem strengthening}

Suppose a local node originally uses residual \(B\), but a new theorem proves
its conclusion without \(B\) or with a weaker precision.  Replacing the local
contract and transfer summary may reduce the global demand closure.  Because
leastness is recomputed, the improvement can propagate across cycles and
remove other resources.  The theorem ledger should record both the local
proof change and the resulting endpoint change.

\subsection{Proof-producing versus proof-checking architectures}

A standalone certificate checker can be useful as an experimental concise
presentation, but DARC's compiler has a different role.  It constructs a term
of the existing \(\Prf\) inductive type.  The local contracts are ordinary
object theorems, and the global schedule uses ordinary structural lemmas.  As a
result:
\begin{itemize}
  \item all existing arithmetic derivations can be reused directly;
  \item the existing axiom audit applies;
  \item open Bridge transport applies;
  \item a client need not trust a second checker or a semantic correspondence
        theorem that has not been integrated.
\end{itemize}
This choice is central to resolving the current repository integration gap.
